\section{已有 E11 实验证据}
\label{sec:e11}

本节整理当前已有实验结果。它们支持本文的局部谱几何机制，并包含 direct activation perturbation 的初步证据；downstream Jacobian 理论仍未被直接验证。

\subsection{更新谱塑形}

E11 结果显示，Muon-like 极分解更新稳定地产生更平、更高秩的 update spectrum。在 noisy mini-batch default setting 的 equal-update control 下，nrUpdate 比值约为 \(2.017\)，stUpdate 比值约为 \(4.765\)。这说明极分解更新确实改变更新矩阵奇异值分布，但这只是机制事实，不是性能声明。

\begin{table}[t]
\centering
\caption{E11 已有局部几何证据摘要。比值均为 Muon-like/Adam 或 flat-polar/GD-spectrum。}
\label{tab:e11-core}
\begin{tabular}{llll}
\toprule
证据项 & 指标 & 数值 & 解释 \\
\midrule
更新谱塑形 & nrUpdate & $2.017\,[1.905,2.136]$ & Muon-like 更新核秩更高 \\
更新谱塑形 & stUpdate & $4.765\,[4.582,4.955]$ & Muon-like 更新稳定秩更高 \\
一阶校准 & $\rho(\Delta L,\langle G,D\rangle)$ & $0.9803\,[0.9787,0.9817]$ & one-step loss 由一阶内积很好解释 \\
Frobenius 预算 & flat/GD & $0.6071\,[0.5846,0.6304]$ & flat/polar 谱分配输给梯度谱分配 \\
Operator 预算 & flat/GD & $1.689\,[1.614,1.768]$ & flat/polar 谱分配胜出 \\
\bottomrule
\end{tabular}
\end{table}


\subsection{一阶损失下降校准}

E11 记录的 \texttt{update\_grad\_inner} 对应本文记号中的
\[
\langle G,D\rangle,
\]
即正的一阶下降 proxy。已有结果显示，one-step 实际损失下降与该 proxy 的 Spearman 相关约为 \(0.9803\)，说明在当前短步长范围内，使用局部一阶模型是合理的。

\subsection{范数几何边界}

谱分配 probe 固定梯度奇异向量，只改变奇异值分配，并分别在 Frobenius 预算与 operator-norm 预算下比较 flat/polar 与 GD-spectrum。结果表明：
\[
\text{Frobenius budget:}\quad \text{flat/GD}=0.6071<1,
\]
但
\[
\text{operator budget:}\quad \text{flat/GD}=1.689>1.
\]
这与附录~\ref{sec:appendix-proofs}中的一阶最优性推导完全一致。

\subsection{神经网络负控}

神经网络 sanity checks 显示，高 update rank 并不自动带来更好一阶进展。Deep MNIST MLP、MNIST Patch surrogate 与 MNIST ConvNet 都保留 Muon-like 更新的高核秩/高稳定秩 signature，但 one-step 对齐比仍显著低于 Adam。见表~\ref{tab:e11-neural-negative}。

\begin{table}[t]
\centering
\caption{神经网络 sanity checks：高秩更新不自动意味着更好一阶进展。}
\label{tab:e11-neural-negative}
\begin{tabular}{llll}
\toprule
任务 & nrUpdate & stUpdate & one-step 对齐比 \\
\midrule
Deep MNIST MLP & $2.262\,[2.175,2.353]$ & $18.31\,[16.26,20.62]$ & $0.5978\,[0.5633,0.6345]$ \\
MNIST Patch + 共享权重 & $4.235\,[3.834,4.677]$ & $12.89\,[11.77,14.12]$ & $0.5222\,[0.4943,0.5516]$ \\
MNIST ConvNet & $3.718\,[3.451,4.006]$ & $12.03\,[11.03,13.12]$ & $0.677\,[0.6321,0.7251]$ \\
\bottomrule
\end{tabular}
\end{table}


这部分证据支持本文的谨慎表述：
\[
\boxed{\text{SpecGrad 改变更新谱几何，但谱扩散是否有利取决于 activation/downstream geometry。}}
\]

\subsection{预测边界尚未完成}

当前边界预测器只能作为描述性分析。已有 audit 显示 in-sample fit 具有探索价值，但 held-out generalization 还不够稳定；部分 held-out setting 还存在 positive/negative label 退化问题。因此，当前不应声称已有一个可迁移的强预测规则。
